\cvsection{Expériences professionnelles}

\begin{cventries}
  \cventry
    {Stagiaire ingénieur de recherche en IA}
    {INRIA}
    {Bordeaux, France}
    {Mars 2026 -- Septembre 2026}
    {
      \begin{cvitems}
        \item {Utilisation de ressources de calcul HPC avec Slurm pour mener des expérimentations et comparer plusieurs modèles de langage.}
        \item {Installation et configuration de modèles, puis déploiement et serving avec vLLM sur une infrastructure HPC équipée de GPU NVIDIA.}
        \item {Conteneurisation de composants avec Docker et mise en place d’une CI/CD sur GitHub, avec versionnement du code sous Git.}
        \item {Gestion d’une application web utilisant des LLM exécutés localement avec LM Studio sur un Mac Studio, avec installation et configuration des modèles.}
        \item {Développement de scripts Python de collecte de métadonnées scientifiques en ligne et transmission des résultats de tâches cron par webhooks.}
        \item {Évaluation de Qwen, Gemma, GLM et de plusieurs quantifications pour étudier les compromis entre qualité et performances selon l’environnement d’exécution.}
        \item {Définition du protocole expérimental, mesure de la latence et du TTFT, analyse des résultats et rédaction en cours d’un article scientifique.}
        \item {Évaluation d’OpenClaw dans un environnement Docker isolé et analyse des risques liés à l’accès autonome aux outils et aux fichiers.}
        \item {Conception et développement en autonomie, avec le suivi de mes encadrants, d’une application et d’un agent IA pour aider les chercheurs à trouver des articles scientifiques.}
        \item {Développement de l’agent en Python sans framework d’orchestration, avec appels d’outils, gestion de la mémoire et du contexte, et intégration de l’API OpenAI.}
        \item {Développement de services Python/FastAPI et d’API REST documentées avec OpenAPI, avec persistance MongoDB et interface React/Next.js.}
      \end{cvitems}
    }

  \cventry
    {Stagiaire ingénieur de recherche en apprentissage par renforcement}
    {Université Karunya}
    {Coimbatore, Inde}
    {Juin 2025 -- Sept. 2025}
    {
      \begin{cvitems}
        \item {Développement d’un modèle d’apprentissage par renforcement et de son environnement d’entraînement pour la navigation autonome d’un robot hospitalier.}
        \item {Déploiement et optimisation du logiciel et du modèle sur Raspberry Pi pour des tests en conditions réelles.}
      \end{cvitems}
    }
\end{cventries}
